\subsection{Critical Reflection}
\label{subsec:kritische-wuerdigung}

The project-specific case study has limited generalizability. The
documentation, acceptance, and local verification originate from the project
context; the developer perspective may influence the selection of evidence.
Literature and official documentation are no substitute for independent
replication. Commit-specific findings remain valid, but their transfer to
other repository states or projects is limited.

\paragraph{Construct validity.}
Code LOC, function and class size, and complexity represent selected structural
properties. They measure neither domain comprehension, cohesion, nor correct
distribution of responsibilities. Import rules do not capture every dynamic,
transitive, or semantic architecture violation. In particular, 900 code LOC
is a generous project-specific maximum, not a universal maintainability
criterion.

\paragraph{Internal validity.}
The evaluation originates within the project. Boundary and architecture
fixtures may be tailored to the implemented rules. Simulated Ruff and ESLint
outputs and the missing negative Rust fixture limit the integration evidence.
Integrated tests for the 901-LOC exception through gate status and for absolute
and excluded exception paths are also absent. Passing tests do not rule out
errors outside the tested cases.

\paragraph{External validity.}
The subject is one profile-based template. Thresholds and architecture mapping
cannot readily be transferred to small starters, large enterprise systems, or
generated codebases. The principles of central policy, declarative variability,
and shared verification entry points are more transferable than the concrete
numeric values.

\paragraph{Reliability.}
Results depend on the commit, tool versions, operating system, package
registries, and native libraries. Divergent master and profile results and the
locally failed Rust/Tauri path demonstrate this dependency. Without coverage
measurement, Compose startup, and successful platform runs, there is no
comprehensive evidence of quality or production readiness.

\paragraph{Long-term effect.}
The study establishes technical detection of defined violations at the time of
observation. Maintenance effort, defect rate, refactoring frequency,
development time, and model compliance were not measured longitudinally. The
quality gate therefore guarantees neither clean code nor a causal improvement
in long-term maintainability.
